\documentclass[11pt,a4paper]{article}
\usepackage[margin=25mm]{geometry}
\usepackage{amsmath,amssymb,amsthm,hyperref}
\newtheorem{lemma}{Lemma}
\newcommand{\gen}[1]{\langle#1\rangle}
\title{Critical odd-degree tails behind exceptional ends}
\author{Human Compression project: L044 continuation}
\date{}
\begin{document}\maketitle
Throughout $G\leq A_n$, $n\geq12$, is transitive with connected
string C-marking, all maximal omissions intransitive, and the
alternating rank bound is available at smaller degrees at least twelve.
This is a bounded reduction, not an exclusion of its final equality.

\begin{lemma}
Suppose the first three generators are an actual supported
six-point $A_5$ marking from source Table~2(2) or (3),
allowing reversal of that displayed triple. If $r>\lfloor(n-1)/2\rfloor$, then
\[
 n=2r\geq18,\qquad H=G_{>2}=A_Z,\qquad
 |Z|=n-5\geq13,\qquad \operatorname{rank}(H)=(|Z|-1)/2.
\]
Here $Z$ is the complement of the first five vertices of the
literal head, and $G_{>3}$ fixes its sixth vertex. All maximal
omissions of $H$ are intransitive on $Z$.
\end{lemma}
\begin{proof}
First settle the orientation literally. In the reversed row~2, its
first two generators are $(1\ 2)(5\ 6)$ and $(2\ 3)(4\ 5)$;
their orbits are the two triples $\{1,2,3\}$ and $\{4,5,6\}$.
They have no fixed point in the head. Any commuting next generator
therefore preserves the head, and every further generator centralizes
the full head. There is no exit, contrary to global transitivity.
Reversal of row~3 is conjugate to the original row by
$(1\ 6)(2\ 5)(3\ 4)$. Thus it suffices to use either displayed
orientation, with no local reversal silently assumed harmless.

Use the exact accepted source-table labels: the first two generators
are $(1\ 2)(3\ 4)$ and $(2\ 3)(4\ 5)$, and the third is either
$(1\ 2)(5\ 6)$ or $(3\ 4)(5\ 6)$. The first two generate the
natural $D_{10}$ on $\{1,\ldots,5\}$ and fix every other point.
Its centralizer on those five points is trivial (centralizing the
five-cycle gives a power of it, and its reflection inverts that
power). Hence $\rho_3$, which commutes with the first two, fixes
those five points. Every $\rho_k$, $k\geq4$, centralizes the entire
head $A_5$, so preserves its unique nontrivial orbit of size six
and fixes it pointwise: the accepted six-point action is
2-transitive and has trivial centralizer. Thus $H$ is supported on
$Z=\Omega\setminus\{1,\ldots,5\}$.

Any $H$-orbit in $Z$ not containing vertex6 would be fixed pointwise
by the head and preserved by all of $G$, contradicting transitivity.
Therefore $H$ is transitive on all of $Z$, and has faithful even
C-marking of rank $d=r-3$. Its endpoint omission fixes vertex6.
A transitive maximal omission of $H$ on $Z$ would make the
corresponding global omission already connect every point moved
by its omitted generator; adjoining that generator could not
join omission orbits. This contradicts transitivity of $G$.

Put $q=|Z|=n-5$. Rank failure gives $d\geq(q-1)/2$.
For $q<12$ the accepted small supported-tail lemma leaves only
five- or six-point groups, whereas $q\geq7$ here. For $q\geq12$
the accepted proper-primitive and anchored-imprimitive bounds
exclude every proper transitive subgroup of $A_Z$ at this rank.
Thus $H=A_Z$. Smaller-degree alternating induction gives
$d\leq\lfloor(q-1)/2\rfloor$. Comparing the two inequalities
forces $q$ odd and $n=2r$. Since $q\geq12$, this gives $q\geq13$
and $n\geq18$ as claimed.
\end{proof}

\begin{lemma}
Suppose $G_0$ has the two orbits $\{a\}$ and $B$. In the same
contradiction range,
\[
 n=2r\geq14,\qquad G_0=A_B,\qquad |B|=n-1\geq13,
 \qquad \operatorname{rank}(G_0)=(|B|-1)/2.
\]
Moreover $G_{>1}$ is intransitive on its support.
\end{lemma}
\begin{proof}
The actual supported tail $G_0=G_{>0}$ has degree $q=n-1$ and
rank $d=r-1\geq(q-1)/2$. Its evenness, IP and intransitive
omissions follow as above. Its external extension supplies the
endpoint anchor used by the complete small and large supported-tail
bounds. Since $q\geq11$, the small survivors of degrees5 and6
are impossible. The proper large cases are likewise excluded, so
smaller-degree alternating induction gives the asserted full odd-degree
alternating equality.
The unique $0$-crossing $ab$ follows from the singleton orbit.
Suppose $K=G_{>1}$ is transitive on its support $X$, of size $x$.
It fixes $a,b$. If $X=\Omega\setminus\{a,b\}$, the complete
accepted singleton-endpoint theorem already gives the global rank
bound. Otherwise $x\leq n-3$ and
$\operatorname{rank}(K)=r-2=(n-4)/2\geq(x-1)/2$.
The actual supported-tail bounds apply again. Their small survivors
have ranks two or three, impossible here since $n\geq14$ gives
$r-2\geq5$. The proper large cases are excluded, so $K=A_X$ and
smaller-degree induction forces $x=n-3$. Its centralizer on $X$
is trivial. Since $\rho_0$ commutes with $K$, it fixes $X$ pointwise.
On the three remaining points it contains $(a\ b)$ and hence is
odd, a contradiction. Thus $K$ is intransitive on its support.
\end{proof}

\paragraph{Limit.}
Neither lemma rules out its final odd-degree alternating marking.
These are precise remaining configurations for the source
Propositions~5.29--5.32 continuation in Cameron--Fernandes--Leemans--Mixer,
\emph{Highest rank of a polytope for $A_n$}, arXiv:1605.09173v1,
physical pages38--44. The literal head data are the accepted
Table~2 transcript from physical page27. The source's end-subgroup
strategy is retained; the actual five-fixed-point tail and its exact
critical range are the present reduction using the complete paid
small, anchored, primitive and induction interfaces. Both literal head orientations have been checked explicitly: reversed
row~2 cannot extend and row~3 has the displayed reversing conjugation.
All full small certificates and predecessor proofs remain paid.
No new computation or formalization is used. This author proof has not yet received distinct review; no whole
L044, original, question or priority credit is claimed.
\end{document}
